\section{Foundational Derivation and Physical Uniqueness}
\label{sec:pr35-foundational-derivation-standard}

The statement that the PR equations are the only solution must be defined over
a mathematical candidate class.  For any fixed number $r$, infinitely many
syntactically different equations can be constructed with $r$ as a root,
for example
\begin{equation}
x-r=0,
\qquad
(x-r)(x^2+1)=0,
\qquad
e^{x-r}-1=0.
\end{equation}
Invertible field redefinitions generate further equivalent descriptions.
Consequently, unrestricted uniqueness among every imaginable equation is not
a meaningful theorem. Let \(\mathfrak A_{\rm PR}\) denote the complete class of operators and equations satisfying the frozen PR postulates, topology, regularity, symmetry, stability, boundary, representation, and declared finite-tier rules.  Let
\(\sim\) identify gauge, basis, coordinate, phase, algebraic, and unit-equivalent
descriptions.  The physical uniqueness target is
\begin{equation}
\boxed{
\operatorname{Sol}(\mathfrak A_{\rm PR})/\!\sim
=
\{[\mathrm{PR}]\}.}
\label{eq:pr35-physical-uniqueness-target}
\end{equation}

The companion derivation-completeness condition requires every published
result $R$ to possess a proved directed path from the frozen base-postulate
set \(\mathcal P_0\):
\begin{equation}
\forall R\in\mathcal R_{\rm PR},
\qquad
\mathcal P_0
\longrightarrow E_1\longrightarrow\cdots\longrightarrow E_k
\longrightarrow R.
\label{eq:pr35-derivation-path}
\end{equation}
The graph must be acyclic, every parent must exist, no empirical diagnostic may be an ancestor of a generated node, and no final theorem output may descend from an unresolved branch selector. Every node is classified as a base postulate, proved derived object, certified numerical object, open selector, finite-tier rule, unit setter, diagnostic, or published output.  Reproducibility of an open selector does not change its classification: a formula can be evaluated exactly and remain underived. The final construction also requires axiom minimality.  Removing each proposed postulate in turn must either make the construction inconsistent or enlarge the physical solution space.  Otherwise that postulate is redundant and must be derived or removed.  This prevents closure from being obtained trivially by promoting every successful published formula to an independent assumption. The PR-I foundational postulate set is frozen. No PR-II or PR-III equation, coefficient, representation, topology, branch rule, or radiative term may be promoted to an additional physical postulate.  Each unresolved selector must instead be derived uniquely from earlier descendants of the frozen postulates, shown to differ only by a quotient equivalence, or removed from the closed theorem scope.  In symbols,
\begin{equation}
\boxed{
\text{open selector}
\not\longrightarrow
\text{new postulate}.}
\label{eq:pr35-no-new-postulate-rule}
\end{equation}
